\documentclass[12pt, a4paper]{extarticle}

% --- UNIVERSAL PREAMBLE BLOCK ---
\usepackage[a4paper, top=2cm, bottom=2cm, left=2.5cm, right=2.5cm]{geometry}
\usepackage{fontspec}

\usepackage[arabic, bidi=basic, provide=*]{babel}

\babelprovide[import, onchar=ids fonts]{arabic}
\babelprovide[import, onchar=ids fonts]{english}

% Set default/Latin font to Sans Serif in the main (rm) slot
\babelfont{rm}{Noto Sans}
% Assign a specific font for Arabic text
\babelfont[arabic]{rm}{Noto Sans Arabic}

% Add because main language is not English
\usepackage{enumitem}
\setlist[itemize]{label=-}
\setlist[enumerate]{label=\arabic*.}
% ---------------------------------

\usepackage{setspace}
\setstretch{1.15}

% حزم إضافة الإطار (البوردر)
\usepackage{tikz}
\usepackage{eso-pic}

\begin{document}

% رسم الإطار الخارجي المزدوج
\AddToShipoutPictureBG{%
  \begin{tikzpicture}[remember picture, overlay]
    \draw[line width=2pt, draw=black] ([xshift=1.2cm,yshift=-1.2cm]current page.north west) rectangle ([xshift=-1.2cm,yshift=1.2cm]current page.south east);
    \draw[line width=0.5pt, draw=black] ([xshift=1.35cm,yshift=-1.35cm]current page.north west) rectangle ([xshift=-1.35cm,yshift=1.35cm]current page.south east);
  \end{tikzpicture}%
}

% --- التاريخ والرقم المرجعي ---
\noindent
\begin{minipage}{0.48\textwidth}
    \textbf{الرقم المرجعي:} .......................
\end{minipage}
\hfill
\begin{minipage}{0.48\textwidth}
\raggedleft
    \textbf{التاريخ:} \hspace{0.2cm} / \hspace{0.2cm} / \hspace{0.5cm} هـ
\end{minipage}

\vspace{0.1cm}
\begin{flushright}
    \textbf{الموافق:} \hspace{0.2cm} / \hspace{0.2cm} / \hspace{0.5cm} م
\end{flushright}

% --- خط فاصل ---
\vspace{0.2cm}
\noindent\rule{\textwidth}{0.8pt}
\vspace{0.3cm}

% --- العنوان ---
\begin{center}
    \textbf{\LARGE طلب إفادة بتطبيق أدوات وتجربة البحث}
\end{center}

\vspace{0.4cm}

% --- المخاطَب ---
\noindent \textbf{السيدة الفاضلة/ مديرة مدرسة فاطمة الزهراء الثانوية بنات}

\noindent \textbf{حفظها الله}

\noindent \textbf{تحية طيبة وبعد،}

\vspace{0.3cm}

% --- الموضوع ---
\noindent \textbf{الموضوع:} طلب إفادة بتطبيق أدوات وتجربة البحث الميدانية.

\vspace{0.3cm}

% --- نص الخطاب ---
\noindent نأمل من سيادتكم التفضل بالموافقة على منحي شهادة (إفادة) تفيد بقيامي بتطبيق \textbf{أدوات وتجربة البحث} الميدانية الخاصة ببحثي بعنوان:

\vspace{0.2cm}
\begin{center}
    \textbf{``التفاعل بين نمط حشد المصادر (تنافسي/ تشاركي) وزمنه ببيئة تعلم إلكتروني وأثره في حل المشكلات والتدفق الذهني لطلاب المرحلة الثانوية''}
\end{center}
\vspace{0.2cm}

\noindent وذلك على عينة من طالبات المدرسة، في الفترة من \textbf{[التاريخ]} إلى \textbf{[التاريخ]}. علماً بأن التطبيق قد شمل:

\begin{enumerate}[nosep, topsep=2pt]
    \item استخدام بيئة التعلم الإلكتروني المقترحة.
    \item تطبيق أدوات القياس (اختبار حل المشكلات، ومقياس التدفق الذهني).
\end{enumerate}

\noindent ونود الإحاطة بأن الباحث قد \textbf{استعان بمعلمة الفصل} في مراحل تطبيق التجربة الميدانية لضمان سير العملية التعليمية وفق المعايير التربوية والأمنية المتبعة، ولتسهيل تفاعل الطالبات مع الأدوات البحثية.

\vspace{0.2cm}
\noindent وهذه الإفادة لتقديمها لجهة الدراسة (الكلية) لاستكمال إجراءات البحث.

\vspace{0.3cm}
\begin{center}
    \textbf{وتفضلوا بقبول فائق الاحترام والتقدير،}
\end{center}

\vspace{0.4cm}

% --- توقيعات الباحث والمشرف ---
\noindent
\begin{minipage}[t]{0.45\textwidth}
    \textbf{المشرف على الرسالة:}\\[0.2cm]
    أ.د/ .....................................\\[0.2cm]
    \textbf{التوقيع:} .............................
\end{minipage}
\hfill
\begin{minipage}[t]{0.48\textwidth}
    \textbf{مقدَّم لسيادتكم من الباحث:}\\[0.2cm]
    إسلام محمد عبد النبي\\[0.2cm]
    \textbf{قسم:} تكنولوجيا التعليم\\[0.2cm]
    \textbf{كلية:} الدراسات العليا للتربية\\[0.2cm]
    \textbf{التوقيع:} .............................
\end{minipage}

\vspace{0.6cm}
\noindent\rule{\textwidth}{0.4pt}
\vspace{0.3cm}

% --- اعتماد مديرة المدرسة ---
\begin{center}
    \textbf{اعتماد مديرة المدرسة}
\end{center}
\vspace{0.2cm}
\noindent
\textbf{الاسم:} .....................................\\[0.2cm]
\textbf{التوقيع:} .....................................\\[0.2cm]
\textbf{الختم:}

\end{document}